\documentclass[10pt,a4paper]{article}
\usepackage[utf8]{inputenc}
\usepackage[T1]{fontenc}
\usepackage{lmodern}
\usepackage[margin=1.8cm,top=2.0cm,bottom=2.0cm]{geometry}
\usepackage{amsmath,amssymb,amsfonts}
\usepackage{graphicx}
\usepackage{xcolor}
\usepackage{tcolorbox}
\tcbuselibrary{skins,breakable}
\usepackage{enumitem}
\usepackage{booktabs}
\usepackage{adjustbox}
\usepackage{microtype}
\usepackage{hyperref}

% Color definitions
\definecolor{brandblue}{HTML}{2C5AA0}
\definecolor{brandgreen}{HTML}{2E7D52}
\definecolor{brandamber}{HTML}{C8881E}
\definecolor{brandred}{HTML}{B23A48}
\definecolor{brandpurple}{HTML}{6A4C93}
\definecolor{brandink}{HTML}{21355E}
\definecolor{bgfill}{HTML}{EAF0F7}

% Custom tcolorbox environments
\newtcolorbox{intuition}[1][Intuition]{
  enhanced, breakable, colback=bgfill, colframe=brandblue,
  fonttitle=\bfseries\color{white}, title={#1},
  boxrule=0.8pt, arc=3pt, left=10pt, right=10pt, top=8pt, bottom=8pt
}

\newtcolorbox{worked}[1][Worked Example]{
  enhanced, breakable, colback=white, colframe=brandgreen,
  fonttitle=\bfseries\color{white}, title={#1},
  boxrule=0.8pt, arc=3pt, left=10pt, right=10pt, top=8pt, bottom=8pt
}

\newtcolorbox{everyday}[1][Everyday Picture]{
  enhanced, breakable, colback=bgfill!50, colframe=brandamber,
  fonttitle=\bfseries\color{white}, title={#1},
  boxrule=0.8pt, arc=3pt, left=10pt, right=10pt, top=8pt, bottom=8pt
}

\newtcolorbox{watchout}[1][Watch Out]{
  enhanced, breakable, colback=white, colframe=brandred,
  fonttitle=\bfseries\color{white}, title={#1},
  boxrule=0.8pt, arc=3pt, left=10pt, right=10pt, top=8pt, bottom=8pt
}

\newtcolorbox{keytake}[1][Key Takeaways]{
  enhanced, breakable, colback=bgfill, colframe=brandpurple,
  fonttitle=\bfseries\color{white}, title={#1},
  boxrule=0.8pt, arc=3pt, left=10pt, right=10pt, top=8pt, bottom=8pt
}

\newenvironment{steps}{
  \begin{enumerate}[label=\textbf{Step \arabic*:}, leftmargin=*]
}{
  \end{enumerate}
}

\newcommand{\housefig}[2]{
  \begin{center}
    \includegraphics[width=0.88\textwidth]{#1}\\
    \small\textcolor{brandink}{\textbf{Figure:} #2}
  \end{center}
}

\setlength{\parindent}{0pt}
\setlength{\parskip}{6pt}

\begin{document}

% Banner Header
\begin{tcolorbox}[enhanced, colback=brandink, colframe=brandink, arc=0pt, left=12pt, right=12pt, top=12pt, bottom=12pt]
  \color{white}
  \large\textbf{ISM Companion Reader} \hfill \textbf{Session 11}\\[4pt]
  \huge\textbf{Analysis of Variance (ANOVA)}\\[4pt]
  \small One-way ANOVA, two-way ANOVA without interaction, Sum of Squares, and F-ratio tests.
\end{tcolorbox}

\vspace{10pt}

\section{Introduction to Analysis of Variance (ANOVA)}

Analysis of Variance (ANOVA) compares means across three or more populations simultaneously. Running multiple two-sample t-tests increases the overall Type I error rate. ANOVA avoids this by testing all means in a single F-ratio test.

\begin{intuition}[Between-Group vs Within-Group Variance]
Imagine testing three teaching methods across classrooms. Total variation in test scores comes from two sources:
\begin{enumerate}
  \item \textbf{Between-Group Variation (Treatment):} Differences caused by using different teaching methods.
  \item \textbf{Within-Group Variation (Error):} Random individual variation among students inside the same classroom.
\end{enumerate}
If between-group variance is significantly larger than within-group noise, the teaching methods produce real differences.
\end{intuition}

\section{One-Way ANOVA Framework}

For $k$ treatment groups and total sample size $n = \sum n_i$, total variance splits into:
\[
\text{SST} = \text{SSTR} + \text{SSE}
\]
where $\text{SST}$ is Total Sum of Squares, $\text{SSTR}$ is Treatment Sum of Squares, and $\text{SSE}$ is Error Sum of Squares.

\housefig{figures/fig_anova_ss.png}{Breakdown of Total Sum of Squares ($\text{SST}$) into Treatment ($\text{SSTR}$) and Error ($\text{SSE}$).}

\textbf{One-Way ANOVA Table:}
\begin{center}
\begin{adjustbox}{max width=\textwidth}
\begin{tabular}{lcccc}
\toprule
\textbf{Source of Variation} & \textbf{Sum of Squares (SS)} & \textbf{Degrees of Freedom (df)} & \textbf{Mean Square (MS)} & \textbf{F-Ratio} \\
\midrule
Treatments (Between) & $\text{SSTR}$ & $k - 1$ & $\text{MSTR} = \frac{\text{SSTR}}{k-1}$ & $F = \frac{\text{MSTR}}{\text{MSE}}$ \\
Error (Within) & $\text{SSE}$ & $n - k$ & $\text{MSE} = \frac{\text{SSE}}{n-k}$ & -- \\
\midrule
Total & $\text{SST}$ & $n - 1$ & -- & -- \\
\bottomrule
\end{tabular}
\end{adjustbox}
\end{center}

\housefig{figures/fig_f_distribution.png}{Fisher F-distribution curve $F(2, 12)$ with critical cutoff $F_{0.05} = 3.89$.}

\begin{worked}[Slide Example: One-Way ANOVA Table]
\textbf{Problem Statement:} A study compares 3 normal populations with 5 samples each ($k = 3, n = 15$). Given Grand Total $G = 179$, $\text{SSTR} = 24.0$, and $\text{SSE} = 72.0$, construct the ANOVA table and test for mean equality at $\alpha = 0.05$.

\begin{steps}
\item \textbf{Hypotheses:} $H_0: \mu_1 = \mu_2 = \mu_3$ vs $H_1:$ at least one mean differs.
\item \textbf{Calculate Degrees of Freedom:}
\[
df_{\text{treatment}} = k - 1 = 2, \quad df_{\text{error}} = n - k = 12, \quad df_{\text{total}} = 14
\]
\item \textbf{Compute Mean Squares:}
\[
\text{MSTR} = \frac{24.0}{2} = 12.0, \quad \text{MSE} = \frac{72.0}{12} = 6.0
\]
\item \textbf{Compute F Test Statistic:}
\[
F_{\text{cal}} = \frac{\text{MSTR}}{\text{MSE}} = \frac{12.0}{6.0} = 2.00
\]
\item \textbf{Decision:} Critical value $F_{0.05}(2, 12) = 3.89$. Since $F_{\text{cal}} = 2.00 < 3.89$, we fail to reject $H_0$. There is no significant difference among the population means.
\end{steps}
\end{worked}

\section{Two-Way ANOVA (Without Interaction)}

In Two-Way ANOVA, data is classified by two factors: Treatments ($k$ groups) and Blocks ($m$ blocks).

\[
\text{SST} = \text{SSTR} + \text{SSBL} + \text{SSE}
\]

\begin{keytake}[Summary Checklist]
\begin{itemize}
  \item \textbf{ANOVA Purpose:} Tests equality of $k \ge 3$ population means simultaneously.
  \item \textbf{F-Statistic:} $F = \text{MSTR} / \text{MSE}$. Large $F$ indicates significant treatment differences.
  \item \textbf{Two-Way ANOVA:} Controls for blocking factor, reducing error variance ($\text{SSE}$).
\end{itemize}
\end{keytake}

\end{document}
